\item \textbf{Chichón en el neumático.}
Mientras viaja detrás de un automóvil a $3.00\text{ m/s}$, nota que una de las llantas del automóvil tiene un pequeño chichón en el borde, como se muestra en la figura P15.16.
\begin{enumerate}
    \item Explique por qué el chichón, desde su punto de vista detrás del automóvil, ejecuta movimiento armónico simple.
    \item Si los radios de los neumáticos del automóvil son de $0.300\text{ m}$, ¿cuál es el periodo de oscilación del chichón?
    \item ¿Qué pasaría si? Usted cuelga un resorte con una constante de resorte $k = 100\text{ N/m}$ desde el espejo retrovisor de su automóvil. ¿Cuál es la masa que debe ser colgada desde este resorte para producir un movimiento armónico simple con el mismo periodo que el chichón en el neumático?
    \item ¿Cuál sería la velocidad máxima de la masa colgante en su automóvil si inicialmente jaló la masa $8.00\text{ cm}$ más allá del equilibrio antes de soltarla?
\end{enumerate}